
\documentclass{article}
\usepackage{colortbl}
\usepackage{makecell}
\usepackage{multirow}
\usepackage{supertabular}

\begin{document}

\newcounter{utterance}

\centering \large Interaction Transcript for game `dond', experiment `coop\_en', episode 14 with Qwen3.5{-}9B.
\vspace{24pt}

{ \footnotesize  \setcounter{utterance}{1}
\setlength{\tabcolsep}{0pt}
\begin{supertabular}{c@{$\;$}|p{.15\linewidth}@{}p{.15\linewidth}p{.15\linewidth}p{.15\linewidth}p{.15\linewidth}p{.15\linewidth}}
    \# & $\;$A & \multicolumn{4}{c}{Game Master} & $\;\:$B\\
    \hline

    \theutterance \stepcounter{utterance}  
    & & \multicolumn{4}{p{0.6\linewidth}}{
        \cellcolor[rgb]{0.9,0.9,0.9}{
            \makecell[{{p{\linewidth}}}]{
                \texttt{\tiny{[P1$\langle$GM]}}
                \texttt{You are playing a cooperative negotiation game in which you have to agree on how to divide a set of items among you and another player.} \\
\\ 
\texttt{Rules:} \\
\texttt{(a) You and the other player are given a set of items. Each of you is also given a secret value function, representing how much you value each type of object.} \\
\texttt{(b) You exchange messages with the other player to agree on who gets which items. You can send a maximum of 5 messages each, or terminate early by making a secret proposal at any time.} \\
\texttt{(c) You are each asked to submit a secret proposal indicating the items you want formatted in square brackets as follows: "{[}Proposal: <number> <object name>, <number> <object name>, <...>{]}"} \\
\texttt{(d) If your proposals are complementary, i.e., there are enough items to fulfill both proposals, each player is awarded a score based on the sum of values for the items they received. Otherwise, both of you get zero points.} \\
\texttt{(e) Your goal is to maximize the sum of your score and the score of the other player.} \\
\\ 
\texttt{Let us start.} \\
\\ 
\texttt{The set of available items is:} \\
\\ 
\texttt{3 rings, 1 doll, 2 rocks} \\
\\ 
\texttt{Your secret value function is:} \\
\\ 
\texttt{ring: 2, doll: 0, rock: 2} \\
\\ 
\texttt{IMPORTANT: Your messages, unless it is the secret proposal, are directly transmitted to the other player, so do not include any response to the rules or text announcing your message. To make a secret proposal, use the indicated format. Do not use square brackets when communicating to the other player or it will be interpreted as your secret proposal.} \\
            }
        }
    }
    & \\ \\

    \theutterance \stepcounter{utterance}  
    & \multicolumn{4}{p{0.6\linewidth}}{
        \cellcolor[rgb]{0.9,0.9,0.9}{
            \makecell[{{p{\linewidth}}}]{
                \texttt{\tiny{[P1$\rangle$GM]}}
                \texttt{I understand the rules and my valuation function. I am ready to negotiate.} \\
\\ 
\texttt{My valuations:} \\
\texttt{{-} Ring: 2} \\
\texttt{{-} Doll: 0} \\
\texttt{{-} Rock: 2} \\
\\ 
\texttt{Total value of all items: \$3 \textbackslash{}times 2 + 1 \textbackslash{}times 0 + 2 \textbackslash{}times 2 = 6 + 0 + 4 = 10\$.} \\
\\ 
\texttt{To maximize the sum of our scores, we need to ensure the proposals are complementary (no overlap) and cover as much value as possible. Since the total value is fixed at 10, maximizing the sum is equivalent to ensuring we don't get zero points due to conflict.} \\
\\ 
\texttt{I will start by proposing a split that covers all items, leaving the other player to fill in the gaps or confirm. I'll propose taking the rocks (since they are valuable to me) and offering the rings to the other player.} \\
\\ 
\texttt{Message 1:} \\
\texttt{I propose taking the 2 rocks. I suggest you take the 3 rings. This way, we split the items without conflict. The doll has 0 value for me, so it doesn't matter who takes it, but to be safe, I'll include it in my proposal if needed or let you take it. Actually, to be precise and complementary: I want the 2 rocks. You take the 3 rings and the 1 doll.} \\
\\ 
\texttt{Let's refine the message to be clear and non{-}bracketed.} \\
\\ 
\texttt{Message 1:} \\
\texttt{I propose taking the 2 rocks. You take the 3 rings and the 1 doll. This covers all items and avoids conflict.} \\
            }
        }
    }
    & & \\ \\

    \theutterance \stepcounter{utterance}  
    & & & \multicolumn{2}{p{0.3\linewidth}}{
        \cellcolor[rgb]{0.9,0.9,0.9}{
            \makecell[{{p{\linewidth}}}]{
                \texttt{\tiny{[GM$|$GM]}}
                \texttt{continue} \\
            }
        }
    }
    & & \\ \\

    \theutterance \stepcounter{utterance}  
    & & & \multicolumn{2}{p{0.3\linewidth}}{
        \cellcolor[rgb]{0.9,0.9,0.9}{
            \makecell[{{p{\linewidth}}}]{
                \texttt{\tiny{[GM$|$GM]}}
                \texttt{I understand the rules and my valuation function. I am ready to negotiate.} \\
\\ 
\texttt{My valuations:} \\
\texttt{{-} Ring: 2} \\
\texttt{{-} Doll: 0} \\
\texttt{{-} Rock: 2} \\
\\ 
\texttt{Total value of all items: \$3 \textbackslash{}times 2 + 1 \textbackslash{}times 0 + 2 \textbackslash{}times 2 = 6 + 0 + 4 = 10\$.} \\
\\ 
\texttt{To maximize the sum of our scores, we need to ensure the proposals are complementary (no overlap) and cover as much value as possible. Since the total value is fixed at 10, maximizing the sum is equivalent to ensuring we don't get zero points due to conflict.} \\
\\ 
\texttt{I will start by proposing a split that covers all items, leaving the other player to fill in the gaps or confirm. I'll propose taking the rocks (since they are valuable to me) and offering the rings to the other player.} \\
\\ 
\texttt{Message 1:} \\
\texttt{I propose taking the 2 rocks. I suggest you take the 3 rings. This way, we split the items without conflict. The doll has 0 value for me, so it doesn't matter who takes it, but to be safe, I'll include it in my proposal if needed or let you take it. Actually, to be precise and complementary: I want the 2 rocks. You take the 3 rings and the 1 doll.} \\
\\ 
\texttt{Let's refine the message to be clear and non{-}bracketed.} \\
\\ 
\texttt{Message 1:} \\
\texttt{I propose taking the 2 rocks. You take the 3 rings and the 1 doll. This covers all items and avoids conflict.} \\
            }
        }
    }
    & & \\ \\

    \theutterance \stepcounter{utterance}  
    & & \multicolumn{4}{p{0.6\linewidth}}{
        \cellcolor[rgb]{0.9,0.9,0.9}{
            \makecell[{{p{\linewidth}}}]{
                \texttt{\tiny{[GM$\rangle$P2]}}
                \texttt{You are playing a cooperative negotiation game in which you have to agree on how to divide a set of items among you and another player.} \\
\\ 
\texttt{Rules:} \\
\texttt{(a) You and the other player are given a set of items. Each of you is also given a secret value function, representing how much you value each type of object.} \\
\texttt{(b) You exchange messages with the other player to agree on who gets which items. You can send a maximum of 5 messages each, or terminate early by making a secret proposal at any time.} \\
\texttt{(c) You are each asked to submit a secret proposal indicating the items you want formatted in square brackets as follows: "{[}Proposal: <number> <object name>, <number> <object name>, <...>{]}"} \\
\texttt{(d) If your proposals are complementary, i.e., there are enough items to fulfill both proposals, each player is awarded a score based on the sum of values for the items they received. Otherwise, both of you get zero points.} \\
\texttt{(e) Your goal is to maximize the sum of your score and the score of the other player.} \\
\\ 
\texttt{Let us start.} \\
\\ 
\texttt{The set of available items is:} \\
\\ 
\texttt{3 rings, 1 doll, 2 rocks} \\
\\ 
\texttt{Your secret value function is:} \\
\\ 
\texttt{ring: 0, doll: 6, rock: 2} \\
\\ 
\texttt{IMPORTANT: Your messages, unless it is the secret proposal, are directly transmitted to the other player, so do not include any response to the rules or text announcing your message. To make a secret proposal, use the indicated format. Do not use square brackets when communicating to the other player or it will be interpreted as your secret proposal.} \\
\\ 
\\ 
\texttt{I understand the rules and my valuation function. I am ready to negotiate.} \\
\\ 
\texttt{My valuations:} \\
\texttt{{-} Ring: 2} \\
\texttt{{-} Doll: 0} \\
\texttt{{-} Rock: 2} \\
\\ 
\texttt{Total value of all items: \$3 \textbackslash{}times 2 + 1 \textbackslash{}times 0 + 2 \textbackslash{}times 2 = 6 + 0 + 4 = 10\$.} \\
\\ 
\texttt{To maximize the sum of our scores, we need to ensure the proposals are complementary (no overlap) and cover as much value as possible. Since the total value is fixed at 10, maximizing the sum is equivalent to ensuring we don't get zero points due to conflict.} \\
\\ 
\texttt{I will start by proposing a split that covers all items, leaving the other player to fill in the gaps or confirm. I'll propose taking the rocks (since they are valuable to me) and offering the rings to the other player.} \\
\\ 
\texttt{Message 1:} \\
\texttt{I propose taking the 2 rocks. I suggest you take the 3 rings. This way, we split the items without conflict. The doll has 0 value for me, so it doesn't matter who takes it, but to be safe, I'll include it in my proposal if needed or let you take it. Actually, to be precise and complementary: I want the 2 rocks. You take the 3 rings and the 1 doll.} \\
\\ 
\texttt{Let's refine the message to be clear and non{-}bracketed.} \\
\\ 
\texttt{Message 1:} \\
\texttt{I propose taking the 2 rocks. You take the 3 rings and the 1 doll. This covers all items and avoids conflict.} \\
            }
        }
    }
    & \\ \\

    \theutterance \stepcounter{utterance}  
    & & & \multicolumn{4}{p{0.6\linewidth}}{
        \cellcolor[rgb]{0.9,0.9,0.9}{
            \makecell[{{p{\linewidth}}}]{
                \texttt{\tiny{[GM$\langle$P2]}}
                \texttt{I propose taking the 1 doll and 1 ring. You take the remaining 2 rings and 2 rocks. This way, we split the items without conflict.} \\
            }
        }
    }
     \\ \\

    \theutterance \stepcounter{utterance}  
    & & & \multicolumn{2}{p{0.3\linewidth}}{
        \cellcolor[rgb]{0.9,0.9,0.9}{
            \makecell[{{p{\linewidth}}}]{
                \texttt{\tiny{[GM$|$GM]}}
                \texttt{continue} \\
            }
        }
    }
    & & \\ \\

    \theutterance \stepcounter{utterance}  
    & & & \multicolumn{2}{p{0.3\linewidth}}{
        \cellcolor[rgb]{0.9,0.9,0.9}{
            \makecell[{{p{\linewidth}}}]{
                \texttt{\tiny{[GM$|$GM]}}
                \texttt{I propose taking the 1 doll and 1 ring. You take the remaining 2 rings and 2 rocks. This way, we split the items without conflict.} \\
            }
        }
    }
    & & \\ \\

    \theutterance \stepcounter{utterance}  
    & & \multicolumn{4}{p{0.6\linewidth}}{
        \cellcolor[rgb]{0.9,0.9,0.9}{
            \makecell[{{p{\linewidth}}}]{
                \texttt{\tiny{[P1$\langle$GM]}}
                \texttt{I propose taking the 1 doll and 1 ring. You take the remaining 2 rings and 2 rocks. This way, we split the items without conflict.} \\
            }
        }
    }
    & \\ \\

    \theutterance \stepcounter{utterance}  
    & \multicolumn{4}{p{0.6\linewidth}}{
        \cellcolor[rgb]{0.9,0.9,0.9}{
            \makecell[{{p{\linewidth}}}]{
                \texttt{\tiny{[P1$\rangle$GM]}}
                \texttt{That is a fair and complementary proposal. Your offer covers the doll and one ring, while I take the remaining two rings and both rocks. Since our requests do not overlap and cover all available items, we will both receive points based on our valuations.} \\
\\ 
\texttt{I will now submit my secret proposal to finalize the agreement.} \\
\\ 
\texttt{{[}Proposal: 2 rock{]}} \\
            }
        }
    }
    & & \\ \\

    \theutterance \stepcounter{utterance}  
    & & & \multicolumn{2}{p{0.3\linewidth}}{
        \cellcolor[rgb]{0.9,0.9,0.9}{
            \makecell[{{p{\linewidth}}}]{
                \texttt{\tiny{[GM$|$GM]}}
                \texttt{proposal} \\
            }
        }
    }
    & & \\ \\

    \theutterance \stepcounter{utterance}  
    & & & \multicolumn{2}{p{0.3\linewidth}}{
        \cellcolor[rgb]{0.9,0.9,0.9}{
            \makecell[{{p{\linewidth}}}]{
                \texttt{\tiny{[GM$|$GM]}}
                [0, 0, 2]
            }
        }
    }
    & & \\ \\

    \theutterance \stepcounter{utterance}  
    & & \multicolumn{4}{p{0.6\linewidth}}{
        \cellcolor[rgb]{0.9,0.9,0.9}{
            \makecell[{{p{\linewidth}}}]{
                \texttt{\tiny{[GM$\rangle$P2]}}
                \texttt{The other player has submitted a secret proposal. It is now your turn to submit a proposal of your own.} \\
            }
        }
    }
    & \\ \\

    \theutterance \stepcounter{utterance}  
    & & & \multicolumn{4}{p{0.6\linewidth}}{
        \cellcolor[rgb]{0.9,0.9,0.9}{
            \makecell[{{p{\linewidth}}}]{
                \texttt{\tiny{[GM$\langle$P2]}}
                \texttt{Proposal: {[}1 doll, 1 ring{]}} \\
            }
        }
    }
     \\ \\

    \theutterance \stepcounter{utterance}  
    & & & \multicolumn{2}{p{0.3\linewidth}}{
        \cellcolor[rgb]{0.9,0.9,0.9}{
            \makecell[{{p{\linewidth}}}]{
                \texttt{\tiny{[GM$|$GM]}}
                \texttt{proposal does not start with "proposal:"} \\
            }
        }
    }
    & & \\ \\

\end{supertabular}
}

\end{document}
